\section{乘积结构；\termTriangleProduct{}}
\label{sec:triangle-product}

先回忆 (Oh99)、(FOh97) 中所述裤子乘积的定义，并把它置于 (FOOO09,
FOOO10a) 及其他较近文献中的一般拉格朗日\personFloer{}理论这一更现代的背景下。

考虑 $(M,\omega)$ 中紧（单调）拉格朗日子流形的三元组 $L_0,L_1,L_2$，以及\personFloer{}复形
\[
 CF^*(L_1,L_0),\qquad CF^*(L_2,L_1),\qquad CF^*(L_2,L_0).
\]
\stfhIndex{I036}{Maodinglagelangriziliuxing}{锚定拉格朗日子流形（anchored Lagrangian submanifold）}
为在 $CF^*$ 上定义\termFilteredTriangleProduct{}，需要使用第 13.7 节引入的锚定拉格朗日子流形
框架。

一劳永逸地固定环境辛流形 $(M,\omega)$ 的底点 $y$。然后在第 13.7 节的意义下，以
$L_i$ 的锚 $\gamma_i$ 配备 $y$，其中
$\gamma_i:[0,1]\to M$ 满足
\[
 \gamma_i(0)=p_i,\qquad\gamma_i(1)=y
\]
（$i=0,\ldots,2$）。再像第 13.7 节那样，把 $\gamma_i$ 与 $\gamma_j$ 拼接，令
\[
 \ell_{ij}=\gamma_i\#\overline\gamma_j
\]
作为 $\Omega(L_i,L_j)$ 中的底道路。类似地，对每个 $i=0,\ldots,k$ 选择拉格朗日
子丛 $\lambda_i\subset\gamma_i^*TM$，并以显然方式定义
\[
 \lambda_{ij}=\lambda_i\#\overline\lambda_j.
\]
\footnote{译者注：当前仅有 $L_0,L_1,L_2$，原书这里仍令 $i=0,\ldots,k$，沿用了
一般多边形的指标范围；译文照录。}

对 $i=0,1,2$，考虑锚定拉格朗日子流形 $(L_i,(\gamma_i,\lambda_i))$。为简化记号，
把给定锚从记号中省略，将相伴滤过\personFloer{}复形简记为 $CF^*(L_i,L_j)$。另记子模
\[
 CF_\lambda^*(L_i,L_j)=\{x\in CF^*(L_i,L_j)\mid\mathfrak v(x)\geq\lambda\}.
\]

现在可以在链层定义\termTriangleProduct{}，依照 (FOOO09)、(Se03a) 的一般记号，记为
\begin{equation}
 \mathfrak m_2:CF^*(L_1,L_0)\otimes CF^*(L_2,L_1)
 \longrightarrow CF^*(L_2,L_0).
\label{eq:triangle-product-chain-level}
\end{equation}
为定义此乘积，遍历所有三元组
\[
 p_1\in L_1\cap L_0,\qquad
 p_2\in L_2\cap L_1,\qquad
 p_0\in L_2\cap L_0,
\]
以及多边形马斯洛夫指标 $\mu(\beta;p_1,p_2;p_0)$，使类
$[w]\in\pi_2(p_1,p_2;p_0)$ 中 $J$-全纯三角形的模空间
$\mathcal M_3(\beta;p_1,p_2;p_0)$ 的相伴解析指标（亦即虚维数）为零，并计数其中
元素。

现在精确定义 $\mathcal M_3(\beta;p_1,p_2;p_0)$。对 $i=0,1,2$（模 $3$ 计），固定
用于定义 $(L_i,L_{i+1})$ 的\personFloer{}边界映射 $\partial$ 的依 $t$ 族
$J_t^{(i)}$。借助（边界）穿孔 $z_i$ 附近给定的解析坐标
$z=\tau+\sqrt{-1}t$，这些数据在
\[
 \dot\Sigma=D^2\setminus\{z_0,z_1,z_2\}
\]
中 $z_0,z_1,z_2$ 的给定邻域上确定一个\termDomainDependentFamily{} $J$。随后把它延拓到整个
$\dot\Sigma$，使得借助给定解析坐标，在各 $z_i$ 处该选择与已经选定的
$J_t^{(i)}$ 相容。

更精确地，设
$\varphi^{(i)}:[0,\infty)\times[0,1]\to\dot\Sigma$ 是各穿孔邻域的给定解析坐标，
并假设其像彼此不交。记
\[
 U^{(i)}\setminus\{z_i\}
 =\operatorname{Image}\bigl(\varphi^{(i)}|_{(1,\infty)\times[0,1]}\bigr),
 \qquad
 (\varphi^{(i)})^{-1}(z)=(\tau(z),t(z))
\]
（$z\in U^{(i)}$）。先令
\begin{equation}
 J(z)=J_{t(z)}^{(i)},\qquad z\in\bigcup_{i=0}^2U^{(i)}.
\label{eq:domain-dependent-j-near-ends}
\end{equation}
利用 $\mathcal J_\omega$ 可缩，把 $J$ 延拓到整个 $\dot\Sigma$，仍以 $J$ 表示延拓。
然后考虑非线性算子 $w\mapsto\overline\partial_Jw$：
\begin{equation}
 \overline\partial_Jw(z)
 :=\frac12\bigl(dw(z)+J(z)dw(z)j(z)\bigr).
\label{eq:domain-dependent-dbar}
\end{equation}
由此可定义如下有限能量模空间。

\begin{definition}
\label{def:holomorphic-triangle-moduli}
定义 $\mathcal M_3(J;L_0,L_1,L_2)$ 为所有对子
$(w,\{z_0,z_1,z_2\})$ 的集合，其中 $\{z_0,z_1,z_2\}\subset\partial D^2$，且
\[
 w:D^2\setminus\{z_0,z_1,z_2\}\longrightarrow M
\]
是满足下列条件的 $J$-全纯映射：
\begin{enumerate}[label=(\arabic*)]
\item $w$ 连续延拓到 $D^2$；
\item $w$ 满足拉格朗日边界条件
\[
 w(\partial_0D^2)\subset L_0,\qquad
 w(\partial_1D^2)\subset L_1,\qquad
 w(\partial_2D^2)\subset L_2,
\]
其中 $\partial_iD^2\subset\partial D^2$ 是 $z_i$ 与 $z_{i+1}$ 之间的弧段
（$i$ 模 $3$）；
\item $w$ 具有有限能量
\[
 E_J(w)=\frac12\int|dw|_J^2=\int w^*\omega<\infty.
\]
\end{enumerate}
\end{definition}

在穿孔 $z_i$ 附近的邻域 $U_i\subset\dot\Sigma$ 上取任意解析坐标
$(\tau,t)\in[0,\infty)\times[0,1]$，对 $w$ 应用命题 14.1.5 所确立的指数收敛，
得到下列引理。

\begin{lemma}
\label{lem:holomorphic-triangle-asymptotics}
对任意 $w\in\mathcal M_3(J_0;L_0,L_1,L_2)$：
\begin{enumerate}[label=(\arabic*)]
\item $w$ 连续延拓到 $D^2$；
\item 存在三元组 $\{p_0,p_1,p_2\}$，其中按模 $3$ 对 $i=1,2,3$ 有
$p_i\in L_{i-1}\cap L_i$，并且 $w(z_i)=p_i$。
\end{enumerate}
特别地，每个这样的 $w$ 自然带有同伦类 $[w]\in\pi_2(p_1,p_2;p_0)$。
\end{lemma}

上述讨论可概括如下。

\begin{proposition}
\label{prop:triangle-moduli-decomposition}
有分解
\[
 \mathcal M_3(J_0;L_0,L_1,L_2)
 =\coprod_{(p_0,p_1,p_2)}\coprod_{\beta\in\pi_2(p_1,p_2;p_0)}
 \mathcal M(\beta;p_1,p_2;p_0),
\]
其中
\[
 \begin{aligned}
 \mathcal M(\beta;p_1,p_2;p_0)=\{(w,(z_0,z_1,z_2))\mid{}&
 \overline\partial_Jw=0,\ [w]=\beta,\\[-2pt]
 &w(z_i)=p_i,\ i=0,1,2\}.
 \end{aligned}
\]
$\widetilde{\mathcal M}_3(J_0;L_0,L_1,L_2)$ 也有类似分解。
\end{proposition}

下面的命题来自第 13.7 节中作用量泛函与锚定拉格朗日子流形相伴的方式。

\begin{proposition}
\label{prop:triangle-area-action-identity}
设
$w:\dot\Sigma=D^2\setminus\{z_0,z_1,z_2\}\to M$ 是任意光滑映射，它在类
$[w]\in\pi_2(p_1,p_2;p_0)$ 中连续延拓到 $D^2$，满足定义
\ref{def:holomorphic-triangle-moduli} 给出的全部条件，但不必 $J$-全纯。以
$\widetilde p:[0,1]\to M$ 表示取值为 $p\in M$ 的常值道路。则
\begin{equation}
 \int w^*\omega
 =\mathcal A_{L_0L_1}([\widetilde p_1,w_0])
 +\mathcal A_{L_1L_2}([\widetilde p_2,w_2])
 -\mathcal A_{L_0L_2}([\widetilde p_0,w_0]).
\label{eq:triangle-area-action-identity}
\end{equation}
\footnote{译者注：原书式 (16.5.14) 第一项与第三项都使用 $w_0$，而中间项使用
$w_2$；依三个输入/输出端通常预期第一项可能为 $w_1$。译文照录。}
\end{proposition}

该恒等式立即给出任意 $w\in\mathcal M_3(\beta;p_1,p_2;p_0)$ 的能量
$E_J(w)$ 之界，因为对此类 $w$ 有 $E_J(w)=\int w^*\omega$。这一能量界结合
$L_0,L_1,L_2$ 的单调性给出下面的紧性。我们会指出
\eqref{eq:domain-dependent-dbar} 中对 $J$ 的规定如何进入证明。这个论证对于较近\personFloer{}理论的进一步发展很重要，尤其是在\personFloer{}理论的哈密顿纤维化背景中，确实必须使用\termDomainDependentFamily{} $J$。

\begin{corollary}
\label{cor:zero-dimensional-triangle-moduli-compact}
若 $\dim\mathcal M_3(\beta;p_1,p_2;p_0)=0$，则
$\mathcal M_3(\beta;p_1,p_2;p_0)$ 紧。
\end{corollary}

\begin{proof}
在 $\dot\Sigma$ 上固定\personKahler{}度量 $h$，它在 $U^{(i)}$（$i=0,1,2$）上是柱形的；
令 $(\tau,t)$ 为 $U^{(i)}$ 的柱形坐标，使 $h=d\tau^2+dt^2$。

设 $w_k$ 是 $\mathcal M_3(\beta;p_1,p_2;p_0)$ 中的序列。先断言：对任意给定的
$z_i$ 的开邻域 $V^{(i)}\subset U^{(i)}$，存在只依赖于这些邻域的 $C>0$，使
\[
 \max_{z\in\dot\Sigma\setminus\bigcup_iV^{(i)}}|dw_k|\leq C.
\]
否则存在子序列 $k_l$ 以及
$z_l\in\dot\Sigma\setminus\bigcup_iV^{(i)}$，使 $|dw(z_l)|\to\infty$，从而在
$\dot\Sigma\setminus\bigcup_iV^{(i)}$ 上产生气泡。但利用 $L_i$ 单调性的维数计算
排除了这种冒泡。因此断言成立。
\footnote{译者注：原书在反证句中从序列 $w_{k_l}$ 改写成未带下标的
$|dw(z_l)|\to\infty$；译文照录。}

因此，由命题 14.2.4，收敛失败唯一可能来自正能量穿过穿孔而损失。用解析柱坐标
$(\tau,t)$ 表述，这对应于一个 $J$-全纯条带从某个穿孔 $z_i$ 分裂出去。确实，若无
能量损失，命题 14.2.4 的同一论证便蕴含该序列的某个子序列收敛。

再由\eqref{eq:domain-dependent-dbar} 对 $J$ 的规定，分裂出的条带是
$\mathcal M(L_i,L_{i+1})$ 中的非常值映射
$u:\mathbb R\times[0,1]\to M$。特别地，其指标至少为 $1$，这与假设
$\dim\mathcal M_3(\beta;p_1,p_2;p_0)=0$ 矛盾。
\end{proof}

所以可以定义
\[
 n(\beta;p_1,p_2;p_0)=\#\mathcal M_3(\beta;p_1,p_2;p_0)
\]
以及
\begin{equation}
 \mathfrak m_2(p_1,p_2)
 =\sum_{p_0\in L_0\cap L_2}\sum_\beta
 n(\beta;p_1,p_2;p_0)T^{\omega(\beta)}e^{\mu(\beta;p_1,p_2;p_0)}p_0,
\label{eq:triangle-product-unlifted}
\end{equation}
并在万有\personNovikov{}环上的复形 $CF^*$ 上线性延拓。严格说来，在引理 13.7.2 的意义
下，此定义应当用于 $[\widetilde p_i,w_i]$，并且
\[
 [w_0]=[w_1]\#[w_2]\#\beta.
\]
但只要同伦类 $[w_i]$ 满足最后这个粘合恒等式，矩阵系数
$n(J_0,\beta;p_1,p_2;p_0)$ 就不依赖于 $[w_i]$，而只依赖于
$\beta\in\pi(p_1,p_2;p_0)$。故上述公式的精确形式为
\begin{equation}
 \mathfrak m_2([\widetilde p_1,w_1],[\widetilde p_2,w_2])
 =\sum_\beta n(\beta;p_1,p_2;p_0)
 [\widetilde p_0,w_1\#w_2\#\beta],
\label{eq:triangle-product-lifted}
\end{equation}
再在万有\personNovikov{}环上的复形 $CF^*$ 上线性延拓。要把它写成
\eqref{eq:triangle-product-unlifted}，必须选择提升
$\widetilde p_i\in\widetilde\Omega_0(L_{i-1},L_i)$，再把 $CF^*(L_{i-1},L_i)$
看成以\personNovikov{}环为系数、在 $L_{i-1}\cap L_i$ 上的自由模，而不是
$\mathbb Q$-向量空间（说明见第 13.9 节）。乘积
\eqref{eq:triangle-product-chain-level} 正是矩阵系数由
$n(J_0,\beta;p_1,p_2;p_0)$ 给出的乘积。

由命题\ref{prop:triangle-area-action-identity} 与 $\mathfrak m_2$ 的定义，还立即得到
映射\eqref{eq:triangle-product-chain-level} 限制为
\[
 \mathfrak m_2:CF_\lambda^*(L_1,L_0)\otimes CF_\mu^*(L_2,L_1)
 \longrightarrow CF_{\lambda+\mu}^*(L_2,L_0).
\]
这是因为，当 $w$ 为 $J$-全纯时，$\int w^*\omega\geq0$；验证是直接的。

\begin{proposition}
\label{prop:triangle-product-chain-map}
乘积 $\mathfrak m_2$ 满足
\begin{equation}
 \partial(\mathfrak m_2(x,y))
 =\mathfrak m_2(\partial(x),y)\pm\mathfrak m_2(x,\partial(y))
\label{eq:triangle-product-leibniz}
\end{equation}
（模去符号），因而在同调上诱导乘积映射
\begin{equation}
 *_F:HF^*(L_1,L_0)\otimes HF_\mu^*(L_2,L_1)
 \longrightarrow HF_{\lambda+\mu}^*(L_2,L_0).
\label{eq:triangle-product-on-homology}
\end{equation}
\footnote{译者注：原书式 (16.5.18) 的第一个因子写作未带滤过下标的
$HF^*(L_1,L_0)$，而式 (16.5.11) 的限制式对应的是 $CF_\lambda^*(L_1,L_0)$；
译文照录。}
\end{proposition}

\begin{proof}
考察满足
\[
 \dim\mathcal M_3(\beta;p_1,p_2;p_0)=1
\]
的四元组 $(\beta;p_1,p_2;p_0)$ 所对应紧化模空间
$\overline{\mathcal M}_3(\beta;p_1,p_2;p_0)$ 的边界。与定理
\ref{thm:boundary-of-index-two-floer-moduli} 的证明类似，$L_0,L_1,L_2$ 的单调性
排除了冒泡。再次利用\eqref{eq:domain-dependent-dbar} 对 $J$ 的规定，得到边界
$\overline{\mathcal M}_3(\beta;p_1,p_2;p_0)$ 由下列三类纤维积组成：
\[
 \begin{gathered}
 \mathcal M_3(\beta_0;p_1,p_2;p'_0)\#\mathcal M(B_0;p'_0,p_0),\\
 \mathcal M(B_1;p_1,p'_1)\#\mathcal M_3(\beta_1;p'_1,p_2;p_0),\\
 \mathcal M(B_2;p_2,p'_2)\#\mathcal M_3(\beta_2;p_1,p'_2;p_0),
 \end{gathered}
\]
其中分别有 $\beta=\beta_0\#B_0$，以及
$\beta=B_i\#\beta_i$（$i=1,2$）；三角形部分维数为 $0$，条带部分维数为 $1$。
应用带边界紧一维流形的分类定理，得到
\[
\begin{aligned}
0={}&\sum_{p'_0\in L_0\cap L_2}\sum_{\beta_0+B_0=\beta}
n(\beta_0;p_1,p_2;p'_0)n(B_0;p'_0,p_0)\\
&\pm\sum_{p'_1\in L_0\cap L_1}\sum_{\beta_1+B_1=\beta}
n(\beta_1;p'_1,p_2;p_0)n(B_1;p_1,p'_1)\\
&\pm\sum_{p'_2\in L_1\cap L_2}\sum_{\beta_2+B_2=\beta}
n(\beta_0;p_1,p'_2;p_0)n(B_2;p_2,p'_2).
\end{aligned}
\]
\footnote{译者注：原书最后一行的求和条件使用 $\beta_2+B_2=\beta$，但对应的
三角形计数写作 $n(\beta_0;p_1,p'_2;p_0)$ 而不是 $n(\beta_2;\ldots)$；译文照录。}
这正是\eqref{eq:triangle-product-leibniz} 的矩阵系数方程（模去符号）。关于符号的
精确考察，参见 (FOOO09) 第 8 章。
\end{proof}
